\documentclass[10pt,a3paper,landscape]{article}
\usepackage[margin=6mm]{geometry}
\usepackage[T1]{fontenc}
\usepackage[utf8]{inputenc}
\usepackage[french]{babel}
\usepackage{lmodern}
\usepackage{microtype}
\makeatletter\def\input@path{{../../../Style/}}\makeatother
\NeedsTeXFormat{LaTeX2e}
\ProvidesPackage{TLPosterCarteMentale}[2026/08/03 Posters et cartes mentales SI]

\RequirePackage{tikz}
\RequirePackage[most]{tcolorbox}
\RequirePackage{xcolor}
\RequirePackage{amsmath,amssymb}
\RequirePackage{graphicx}
\usetikzlibrary{positioning,arrows.meta,calc,shapes.geometric,mindmap,backgrounds,fit}

\definecolor{TLPosterBlue}{RGB}{24,86,141}
\definecolor{TLPosterLight}{RGB}{234,243,250}
\definecolor{TLPosterAccent}{RGB}{232,126,34}
\definecolor{TLPosterGreen}{RGB}{52,152,102}
\definecolor{TLPosterRed}{RGB}{192,57,43}
\definecolor{TLPosterGray}{RGB}{80,90,100}

\newcommand{\TLPosterTitre}[2]{%
  \begin{tcolorbox}[enhanced,boxrule=0pt,arc=0pt,colback=TLPosterBlue,left=6mm,right=6mm,top=4mm,bottom=4mm]
    {\color{white}\bfseries\LARGE #1}\hfill{\color{white}\large #2}
  \end{tcolorbox}%
}

\newtcolorbox{TLPosterBloc}[2][]{
  enhanced,breakable,
  colback=TLPosterLight,
  colframe=TLPosterBlue,
  colbacktitle=TLPosterBlue,
  coltitle=white,
  fonttitle=\bfseries,
  title={#2},
  boxrule=0.8pt,
  arc=2mm,
  left=3mm,right=3mm,top=2mm,bottom=2mm,
  #1
}

\newtcolorbox{TLPosterFormule}[1][]{
  enhanced,
  colback=white,
  colframe=TLPosterAccent,
  boxrule=1pt,
  arc=2mm,
  fontupper=\bfseries,
  halign=center,
  #1
}

\newcommand{\TLPosterMethode}[1]{%
  \begin{tcolorbox}[enhanced,colback=TLPosterGreen!8,colframe=TLPosterGreen,title=Méthode,fonttitle=\bfseries]
  #1
  \end{tcolorbox}%
}

\newcommand{\TLPosterAttention}[1]{%
  \begin{tcolorbox}[enhanced,colback=TLPosterRed!6,colframe=TLPosterRed,title=Point de vigilance,fonttitle=\bfseries]
  #1
  \end{tcolorbox}%
}

\tikzset{
  TLmindroot/.style={concept,concept color=TLPosterBlue,text=white,font=\bfseries\large,minimum size=34mm},
  TLmindlevel1/.style={level 1 concept/.append style={font=\bfseries,minimum size=23mm,sibling angle=45}},
  TLmindlevel2/.style={level 2 concept/.append style={font=\small,minimum size=17mm}},
  TLarrow/.style={-{Stealth[length=3mm]},thick,draw=TLPosterBlue},
  TLbox/.style={rounded corners=2mm,draw=TLPosterBlue,fill=TLPosterLight,align=center,inner sep=3mm}
}

\newenvironment{TLCarteMentale}[1]{%
  \begin{tikzpicture}[mindmap,grow cyclic,concept color=TLPosterBlue,every node/.style={concept},TLmindlevel1,TLmindlevel2]
  \node[TLmindroot]{#1}
}{;
  \end{tikzpicture}
}

\newcommand{\TLBranche}[3][]{child[concept color=#2]{node[#1]{#3}}}

\endinput

\pagestyle{empty}
\renewcommand{\familydefault}{\sfdefault}
\begin{document}
\begin{tikzpicture}
\TLMapBase{Actions mécaniques}{Objectif $\rightarrow$ isolement $\rightarrow$ BAME $\rightarrow$ bon principe $\rightarrow$ vérification}
\TLMapBranchTL{À connaître}{\textbullet\ torseur : résultante + moment ; transport $\vec M_B=\vec M_A+\overrightarrow{BA}\wedge\vec R$\\[1mm]\textbullet\ glisseur, couple, action-réaction\\[1mm]\textbullet\ torseurs transmissibles par les liaisons\\[1mm]\textbullet\ PFS, TEC, PFD\\[1mm]\textbullet\ Coulomb : adhérence / glissement\\[1mm]\textbullet\ centre d'inertie, inertie, Huygens}
\TLMapBranchTR{Choisir la méthode}{\textbf{Équilibre :} PFS\\[1mm]\textbf{Couple moteur / accélération / 1 ddl :} TEC souvent plus rapide\\[1mm]\textbf{Réactions de liaison + mouvement :} PFD\\[1mm]\textbf{Contact rugueux :} Coulomb + vérification du mode\\[1mm]\textbf{Inertie composée :} symétries + Huygens + somme}
\TLMapBranchML{Méthode statique}{1. Graphe de structure.\\[1mm]2. Choisir l'isolement.\\[1mm]3. Dresser le BAME.\\[1mm]4. Réduire les torseurs au point utile.\\[1mm]5. Appliquer $\sum\vec R=\vec0$, $\sum\vec M_A=\vec0$.\\[1mm]6. Projeter seulement les équations utiles.\\[1mm]\textbf{Raccourcis :} 2 actions $\Rightarrow$ colinéaires ; 3 glisseurs $\Rightarrow$ concurrence.}
\TLMapBranchMR{Frottement / contact}{\textbf{Adhérence supposée :} $|T|\le fN$.\\[1mm]\textbf{Limite / glissement :} $|T|=fN$.\\[1mm]Direction de $\vec T$ opposée au glissement (ou au glissement imminent).\\[1mm]\textbf{Méthode :} choisir le mode $\rightarrow$ écrire les équations $\rightarrow$ résoudre $\rightarrow$ vérifier a posteriori l'inégalité de Coulomb.}
\TLMapBranchBL{Inertie / énergie}{\textbf{Centre d'inertie :} décomposer $\rightarrow$ barycentre $\rightarrow$ vérifier symétries.\\[1mm]\textbf{Matrice :} base principale au centre $G$ $\rightarrow$ changement de base $\rightarrow$ Huygens $\rightarrow$ addition/soustraction.\\[1mm]\textbf{TEC :} paramétrer $\rightarrow$ calculer $E_c$ $\rightarrow$ puissances ext./int. $\rightarrow$ $dE_c/dt=\sum\mathcal P$.}
\TLMapBranchBR{PFD / vérifications}{\textbf{PFD :} cinématique $\rightarrow$ torseur dynamique $\rightarrow$ torseurs extérieurs $\rightarrow$ réduction au bon point $\rightarrow$ projections.\\[1mm]\textbf{Vérifier :} équations/inconnues ; signes ; unités ; sens des efforts ; condition de frottement ; ordre de grandeur ; cohérence avec le mouvement.}
\TLMapRibbon{Ne pas choisir PFS, TEC ou PFD par habitude : choisir l'outil qui donne directement la grandeur recherchée avec le minimum d'inconnues.}
\end{tikzpicture}
\end{document}
